\documentclass[10pt,a4paper]{amsart}
\usepackage[margin=19mm]{geometry}
\usepackage{amsmath,amssymb,mathtools}
\usepackage[scr=rsfs,cal=euler]{mathalfa}
\usepackage{xfrac}
\usepackage[unicode,hidelinks,bookmarks=false]{hyperref}
\newcommand{\ie}{i.e.\ }
\newcommand{\cf}{cf.\ }
\newcommand{\resp}{respectively\ }
\newcommand{\Subsection}{\subsection}
\providecommand{\nobreakdash}{\nobreak\hbox{-}}
\setlength{\emergencystretch}{2em}
\multlinegap0pt
\usepackage{mathtools}
\usepackage{amssymb}
\usepackage{enumerate}
\newtheorem{lemma}{Lemma}
\title{Time derivative estimates for normalized parabolic $p$-Laplace equations}
\author{Se-Chan Lee, Hyungsung Yun}
\date{Source: arXiv:2609.15501v1 (CC BY 4.0)}
\begin{document}
\maketitle

\noindent\emph{Source and licence.} Time derivative estimates for normalized parabolic $p$-Laplace equations. Version arXiv:2609.15501v1 (CC BY 4.0), distributed by arXiv under the Creative Commons Attribution 4.0 International licence, http://creativecommons.org/licenses/by/4.0/, as recorded in the arXiv licence metadata for this exact version and shown on the abstract page https://arxiv.org/abs/2609.15501v1. Full licence notice: \texttt{licenses/arxiv-2609.15501-CC-BY-4.0.txt}.\par\smallskip

\noindent\textit{Editorial scope.} Counted unit: the article's lemma lem:logarithmic (source0.tex lines 525--532). The article's complete original proof of it is delivered with it (source0.tex lines 534--646, one \texttt{proof} environment of 113 lines). No mathematics is written, repaired or re-derived: the statement and the proof are the article's own lines, in the article's own notation, and no retained line is substituted or re-flowed.  Retained uncounted as the source-local context the proof needs: lines 280--283; lines 286--291; lines 303--305.  Nothing was omitted from any retained span, so the package carries no omission record.  No retained line contains a citation, so the wrapper carries no bibliography and no bibliography entry.  No retained line contains a figure environment, an image-inclusion command or drawing code, so no image is delivered and the package declares no asset dependency.  Editorial formatting only, in addition: an AMS \texttt{amsart} preamble that restates the article's own declarations and macros used by the retained lines, namely the environments \texttt{lemma} on separate counters, together with the standard packages those lines need.  The retained environments print in this document as Lemma 1 in order of appearance, whereas the article prints the same items with the numbering of its own full document.  The complete native source of the article as distributed in the arXiv e-print for this exact version is bundled unchanged as \texttt{sources/210423/source0.tex}.
\begin{equation}\label{eq:matrix}
 a^{ij}_\varepsilon(q)\coloneqq \delta_{ij}+(p-2)\frac{q_iq_j}{\varepsilon^2+|q|^2}
 \quad \text{for $\varepsilon \in (0, 1]$ and $q \in \mathbb{R}^n$}.
\end{equation}

\begin{equation}\label{eq:ellipticity}
 \lambda_p|\xi|^2
 \leq a_\varepsilon^{ij}(q)\xi_i\xi_j
 \leq\Lambda_p|\xi|^2
 \quad \text{for every $q,\xi\in\mathbb{R}^n$}.
\end{equation}

\begin{equation}\label{eq:gradient}
    \|u^\varepsilon\|_{C^{1,\alpha}(\overline{Q}_{3/4})}\leq C \quad \text{and}  \quad \|u^\varepsilon\|_{C^\alpha(\overline{Q}_{7/8})}\leq C.
\end{equation}

\begin{lemma}[Logarithmic time derivative estimates]\label{lem:logarithmic}
Let $1/2\leq s<s+d\leq2/3$ and $\varepsilon \in (0,1]$. Then we have
\begin{equation}\label{eq:logarithmic}
 \|u_t^\varepsilon\|_{L^\infty(Q_s)}
 \leq Cd^{-2}\log\frac e\varepsilon,
\end{equation}
where $C>0$ depends only on $n$ and $p$.
\end{lemma}

\begin{proof}
For notational convenience, we omit $\varepsilon$ from $u^\varepsilon$ and $a_\varepsilon^{ij}$. We consider the linearized operator
\begin{equation*}
 Lw\coloneqq a^{ij}(Du)w_{ij}+b^l w_l \quad \text{for $b^l\coloneqq a^{ij}_{q_l}(Du)u_{ij}$}.
\end{equation*}
We also write $b=(b^1, \ldots, b^n)$. Differentiation of the equation $u_t=a^{ij}(Du)u_{ij}$ with respect to the variables $t$ and $x_k$ gives, respectively:
\begin{equation}\label{eq:linearized}
 (\partial_t-L)u_t=0
 \quad \text{and} \quad
 (\partial_t-L)u_k=0.
\end{equation}
Moreover, it follows from \eqref{eq:matrix} and \eqref{eq:ellipticity} that
\begin{equation}\label{eq:drift}
    \begin{aligned}
    |u_t|&=|a^{ij}(Du)u_{ij}|\leq C\|D^2u\|,\\
    |b^l|&=|p-2|\left|\left[\frac{\delta_{il}u_j+\delta_{jl}u_i}{\varepsilon^2+|Du|^2}-\frac{2u_iu_ju_l}{(\varepsilon^2+|Du|^2)^2}\right] u_{ij} \right|\leq\frac{C\|D^2u\|}{(\varepsilon^2+|Du|^2)^{1/2}}.
    \end{aligned}
\end{equation}

For the smooth convex function
\begin{equation*}
\Psi_\varepsilon(q)\coloneqq\rho\log\frac\rho\varepsilon \quad \text{for $\rho=(\varepsilon^2+|q|^2)^{1/2}$},
\end{equation*}
it follows from a direct calculation that
\begin{equation*}
    D^2_q\Psi_{\varepsilon}(q)=\frac{1+\log(\rho/\varepsilon)}{\rho}I-\frac{\log(\rho/\varepsilon)}{\rho^3}q \otimes q \geq \frac{1}{\rho}I
\end{equation*}
and that
\begin{equation*}
    \begin{aligned}
        (\partial_t-L)\Psi_\varepsilon(Du)&=(\Psi_{\varepsilon})_{q_k}(Du) (\partial_t u_k-a^{ij}(Du)u_{kij}-b^lu_{kl}) \\
        &\quad-a^{ij}(Du)(\Psi_\varepsilon)_{q_kq_l}(Du)u_{ki}u_{lj}\\
        &=-a^{ij}(Du)(\Psi_\varepsilon)_{q_kq_l}(Du)u_{ki}u_{lj}.
    \end{aligned}
\end{equation*}
Since $(a^{ij}) \geq \lambda_pI$ and $D^2\Psi_{\varepsilon} \geq \rho^{-1}I$, we have
\begin{equation*}   
        K \coloneqq a^{ij}(Du)(\Psi_\varepsilon)_{q_kq_l}(Du)u_{ki}u_{lj}\geq\frac{\lambda_p\|D^2u\|^2}{(\varepsilon^2+|Du|^2)^{1/2}},
\end{equation*}
and so \eqref{eq:drift} yields that
\begin{equation}\label{eq:absorb}
 |u_t||b|\leq CK.
\end{equation}
On the other hand, the gradient bound \eqref{eq:gradient} and $\varepsilon \in (0,1]$ imply
\begin{equation*}
     0\leq\Psi_\varepsilon(Du)\leq C\log\frac e\varepsilon \quad \text{and} \quad 
     K\geq c\|D^2u\|^2 \quad\text{in $Q_{3/4}$}.
\end{equation*}

We now choose a smooth cut-off function $\eta$ on $\overline{Q}_{s+d}$ such that $0\leq\eta\leq1$, $\eta=1$ on $Q_s$, and $\eta=0$ on $\partial_pQ_{s+d}$, with
\begin{equation*}
    |\eta_t|+\|D^2\eta\|\leq Cd^{-2}, \quad
    |D\eta|\leq Cd^{-1}\quad\text{and}\quad
    |D\eta|^2\leq Cd^{-2}\eta.
\end{equation*}
We next let
\begin{equation*}
 A\coloneqq\|\eta u_t\|_{L^\infty(Q_{s+d})}.
\end{equation*}
There is nothing to prove if $A=0$, so we assume that $A>0$. It follows from \eqref{eq:linearized} that
\begin{equation*}
    \begin{aligned}
        (\partial_t-L)(\eta^2u_t^2)=&2\eta(\eta_t-a^{ij}\eta_{ij}-b^l\eta_l)u_t^2
 -2a^{ij}\eta_i\eta_j u_t^2\\
 &-2\eta^2a^{ij}u_{ti}u_{tj}-8\eta u_ta^{ij}\eta_i u_{tj}\\
 &\leq 2\eta(|\eta_t|+|a^{ij}\eta_{ij}|+|b||D\eta|)u_t^2
 -\eta^2a^{ij}u_{ti}u_{tj}+8|\eta u_ta^{ij}\eta_i u_{tj}|.
    \end{aligned}
\end{equation*}
By Young's inequality together with $(a^{ij}) \leq \Lambda_pI$, we have
\begin{equation*}
 8|\eta u_ta^{ij}\eta_i u_{tj}| \leq 8|\eta u_t|(a^{ij}\eta_i\eta_j)^{1/2}(a^{ij}u_{ti}u_{tj})^{1/2}\leq \eta^2a^{ij}u_{ti}u_{tj}+C|D\eta|^2u_t^2.
\end{equation*}
Moreover, we observe that
\begin{itemize}
    \item since $|\eta_t|+|a^{ij}\eta_{ij}| \leq Cd^{-2}$, we have
    \begin{equation*}
        2\eta(|\eta_t|+|a^{ij}\eta_{ij}|)u_t^2 \leq Cd^{-2} \eta u_t^2 \leq Cd^{-2}A|u_t|;
    \end{equation*}

    \item since $|D\eta| \leq Cd^{-1}$, we have
    \begin{equation*}
        2\eta |b||D\eta|u_t^2 \leq Cd^{-1} \eta u_t^2 |b| \leq Cd^{-1} A |u_t| |b|;
    \end{equation*}

    \item since $|D\eta|^2 \leq Cd^{-2}\eta$, we have
    \begin{equation*}
        C|D\eta|^2u_t^2 \leq Cd^{-2}\eta u_t^2 \leq Cd^{-2}A|u_t|.
    \end{equation*}
\end{itemize}
A combination of all estimates gives that
\begin{equation}\label{eq:product}
 (\partial_t-L)(\eta^2u_t^2)
 \leq Cd^{-2}A\|D^2u\|+Cd^{-1}A|u_t||b|.
\end{equation}
Thus, \eqref{eq:absorb} and $\|D^2u\|\leq (\|D^2u\|^2+1)/2\leq CK+1/2$ give
\begin{equation*}
    (\partial_t-L)(\eta^2u_t^2) \leq Cd^{-2}A(K+1)=-Cd^{-2}A (\partial_t-L)(\Psi_{\varepsilon}(Du)-t).
\end{equation*}
We then choose a sufficiently large constant $C_1>0$, depending only on $n$ and $p$, so that the auxiliary function
\begin{equation*}
     v\coloneqq\eta^2u_t^2 + C_1d^{-2}A(\Psi_\varepsilon(Du)-t)
\end{equation*}
satisfies $(\partial_t-L)v<0$ in $Q_{s+d}$. Since
\begin{equation*}
     v=C_1d^{-2}A\Psi_\varepsilon(Du)-C_1d^{-2}At\leq Cd^{-2}A\log\frac e\varepsilon \quad \text{on $\partial_pQ_{s+d}$},
\end{equation*}
the maximum principle gives the same bound for $v$ in $Q_{s+d}$. Finally, we arrive at
\begin{equation*}
    A^2=\|\eta u_t\|_{L^{\infty}(Q_{s+d})}^2 \leq Cd^{-2}A\log\frac e\varepsilon,
\end{equation*}
which finishes the proof by recalling that $\|u_t\|_{L^{\infty}(Q_s)} \leq A$.
\end{proof}

\end{document}
